\section{Baseline Model}

\begin{frame}
  \frametitle{Players I: Consumers and Competitive Benchmark}
  \footnotesize

  \begin{block}{Consumers}
    \begin{itemize}
      % \item Unit demand: each consumer buys at most one product.
      \item Outside option of not buying anything (e.g., postpone purchase or buy another product): $\nu_o \sim G$.\footnote{$G$ is log-concave and continuously differentiable on $[0,\bar{\nu}_o]$, with density $g$.}
      \item Platform transaction benefit: $b>0$ (finding products, one-click payment, delivery tracking, easy returns).
      \item Choice rule: buy the option with the highest net utility.
    \end{itemize}
  \end{block}

  \begin{block}{Fringe sellers}
    \begin{itemize}
      \item At least two identical sellers.
      \item Sell a basic product with value $\nu$.
      \item Consumers already know them and can buy directly.
      \item Bertrand competition gives the price benchmark.\footnote{All marginal costs are normalized to zero.}
    \end{itemize}
  \end{block}
\end{frame}

\begin{frame}
  \frametitle{Players II: Platform and Innovative Seller}
  \small

  \begin{block}{Innovative (superior) seller $S$}
    \begin{itemize}
      \item Benefits from innovation; discovered only if listed on $M$.
      \item Better product value: $\nu+\Delta>\nu>0$.
      \item Innovation: chooses quality investment $\Delta$.
    \end{itemize}
  \end{block}

  \begin{block}{Platform $M$}
    \begin{itemize}
      \item Chooses marketplace, seller, or dual mode.
      \item Seller: sells its own offering, valued at $\nu+\sigma$.
      \item Marketplace: sets commission $\tau \geq 0$ for third-party sellers.
    \end{itemize}
  \end{block}

  
\end{frame}

\begin{frame}
  \frametitle{Values and Parameters}
  \small
  Product values are:

  \begin{center}
    \begin{tabular}{lll}
      \toprule
      Product & Consumer value & Meaning \\
      \midrule
      Fringe sellers & $\nu$ & Basic product \\
      Innovative seller $S$ & $\nu+\Delta$ & Better product \\
      Platform $M$ & $\nu+\sigma$ & Platform's own product \\
      \bottomrule
    \end{tabular}
  \end{center}

  \begin{itemize}
    \item $b>0$: convenience benefit from buying through the platform.
    \item $\nu_o \sim G$: outside option.
    \item $\tau$: platform commission.
  \end{itemize}

  \begin{block}{Intuition}
    $\sigma>0$ means the platform is a better seller than fringe sellers; $\sigma<0$ means the opposite.
  \end{block}
\end{frame}

\begin{frame}
  \frametitle{Innovation by Seller $S$}
  \small
  Seller $S$ chooses innovation $\Delta \geq \Delta^l$ and pays fixed cost $K(\Delta)$.

  \begin{equation*}
    K'(\bar{\Delta}) = G(\nu+\sigma+b)
  \end{equation*}

  \begin{block}{Meaning of $\bar{\Delta}$}
    $\bar{\Delta}$ is the high innovation level that can arise when the market gives strong rewards to innovation.
  \end{block}

  \begin{alertblock}{Intuition}
    The left side is marginal cost of innovation. The right side is the demand gain from better quality.
  \end{alertblock}
\end{frame}

\begin{frame}
  \frametitle{Product Discovery}
  \small
  Consumers do not know about $S$ at the start.

  \begin{itemize}
    \item Consumers know the fringe sellers.
    \item Consumers may know the platform's own product.
    \item Consumers discover $S$ only if $S$ is available on the platform.
  \end{itemize}

  \begin{block}{Why this matters}
    The platform is not only a sales channel. It is also a discovery channel.
  \end{block}

  \begin{alertblock}{Intuition}
    If a policy changes the platform's mode, it may also change whether consumers find $S$.
  \end{alertblock}
\end{frame}

\begin{frame}
  \frametitle{Three Platform Modes}
  \small
  \begin{center}
    \begin{tabular}{lll}
      \toprule
      Mode & What $M$ does & Key feature \\
      \midrule
      Marketplace & Hosts third-party sellers & Earns commission $\tau$ \\
      Seller & Sells its own product & No marketplace for $S$ \\
      Dual & Hosts sellers and sells own product & Both roles at once \\
      \bottomrule
    \end{tabular}
  \end{center}

  \begin{block}{Intuition}
    The policy debate is about whether the third mode should be allowed.
  \end{block}
\end{frame}

\begin{frame}
  \frametitle{Timing}
  \small
  The game is solved by backward induction.

  \begin{enumerate}
    \item $M$ chooses its mode and sets commission $\tau$ if there is a marketplace.
    \item $S$ decides whether to join the platform and chooses $\Delta$.
    \item Sellers set prices.
    \item Consumers choose what to buy.
  \end{enumerate}

  \begin{block}{Equilibrium concept}
    Subgame Perfect Nash Equilibrium.
  \end{block}
\end{frame}
